\section{P1: transport polynomial identities into internal arithmetic}

\subsection{The proof obligation left by a semiring evaluator}

The maintained development has three relevant layers.
\src{Representation}{\code{Expr}} stores integer arithmetic syntax and its
ordinary integer polynomial.
\src{NaturalExpression}{\code{splitNatural}} translates a signed expression
into two natural-coefficient expressions.
\src{GatesSemiring}{\code{SemiringZero}} interprets vanishing as equality of
their natural values.

If $e$ has split $(e^+,e^-)$, write
\[
 \SZ_R(e,v)\quad:\Longleftrightarrow\quad
 [\![ e^+]\!]_R(v)=[\![ e^-]\!]_R(v).
\]
Subtraction need not exist on $R$. An existing semiring evaluator does not
automatically permit replacing $e$ by another expression $f$ with the same
integer polynomial: their positive and negative syntax pieces can look quite
different. Repeatedly unfolding \code{splitNatural} exposes this implementation
to every client.

The proposed interface is
\[
 \boxed{\quad e.\mathrm{toPolynomial}=f.\mathrm{toPolynomial}
 \ \Longrightarrow\
 \bigl(\SZ_R(e,v)\leftrightarrow\SZ_R(f,v)\bigr).\quad}
\]
It holds for every additively cancellative commutative semiring $R$.
A proof can establish a finite algebraic normalization in an ordinary integer
polynomial ring and then reuse it on any appropriate internal-natural carrier.
The stored expression itself remains unchanged.

\subsection{Proof by a natural cross-sum identity}

Let $P_e,N_e,P_f,N_f\in\N[\mathbf X]$ be the natural-coefficient polynomials
of the four split expressions. Let
\[
 \iota:\N[\mathbf X]\longrightarrow\Z[\mathbf X]
\]
be the homomorphism induced by the natural-to-integer coefficient map.

One structural induction proves
\[
 \iota(P_e)-\iota(N_e)=e.\mathrm{toPolynomial}.
 \tag{1}
\]
The constant case is the integer identity
$a.\mathrm{toNat}-(-a).\mathrm{toNat}=a$.
For multiplication, the positive split is $P_eP_f+N_eN_f$ and the negative
split is $P_eN_f+N_eP_f$, so the invariant is the product-of-differences
identity. No semantic model occurs in this induction.

If $e$ and $f$ have equal integer polynomials, (1) gives
\[
 \iota(P_e)-\iota(N_e)=\iota(P_f)-\iota(N_f),
\]
hence
\[
 \iota(P_e+N_f)=\iota(P_f+N_e).
\]
Injectivity of the coefficient map, and therefore its multivariate-polynomial
map, reflects this to
\[
             P_e+N_f=P_f+N_e\quad\text{in }\N[\mathbf X].
 \tag{2}
\]
Evaluate (2) through the homomorphism induced by \code{Nat.castRingHom}~\code{R}:
\[
             p_e+n_f=p_f+n_e.
 \tag{3}
\]
If $p_e=n_e$, substitute and cancel a common additive term to obtain
$p_f=n_f$. The converse is symmetric.

The proof does not embed the semantic carrier in a ring or construct its
Grothendieck completion. It translates a finite coefficient identity once
and retains the original semiring semantics.

\subsection{Checked Lean interface}

The archive's \code{PolynomialZeroTransport.lean} proves four facts about
the actual repository syntax:

\begin{Lean}
map_splitNatural_sub
splitNatural_cross_sum
splitNatural_eval_cross_sum
splitNatural_zero_iff_of_polynomial_eq
\end{Lean}

Its main theorem is:

\begin{Lean}
theorem splitNatural_zero_iff_of_polynomial_eq
    [CommSemiring R] [IsCancelAdd R]
    (e f : Expr σ)
    (h : e.toPolynomial = f.toPolynomial) (v : σ → R) :
    (e.splitNatural.1.evalSemiring v =
       e.splitNatural.2.evalSemiring v) ↔
    (f.splitNatural.1.evalSemiring v =
       f.splitNatural.2.evalSemiring v)
\end{Lean}

A production wrapper using the existing \code{SemiringZero} name is
definitionally the same statement. The core belongs below the gate and
object-model layers to avoid an import cycle. Its checked file imports
\code{NaturalExpressionSemiring}, not the deep object-model closure.

\subsection{Concrete client: the conditional selection row}

The literal row in
\src{ObjectConditional}{\code{literal_selection}} is
\[
       e=o-\bigl(bz+(1-b)p\bigr).
\]
The object-level goal is the balanced equality
\[
       o+bp=bz+p.
 \tag{4}
\]
Normalize to
\[
       f=(o+bp)-(bz+p).
\]
The equality of their ordinary integer polynomials follows by \code{ring}.
P1 transfers the zero predicate from $e$ to $f$; the latter has immediate
split semantics (4).

The checked client leaves all four variable indices abstract and then
instantiates the exact slots used by \code{ObjectConditional}: output $0$,
bit $4$, zero-branch $2$, and positive-branch $3$, with an empty register
block. This checks the actual row shape, not a numerical analogue.

This row already has a short direct proof. It demonstrates the interface,
not a net saving by itself. The useful adoption experiment is a group of
algebraically normalized rows in the zero-test, conditional, pairing, and
division syntax layers. Purely structural gates can retain their current
cheap specialized lemmas.

\subsection{The hypotheses are substantive}

Additive cancellation cannot be removed. In the Boolean semiring, with
addition OR and multiplication AND, consider
\[
              e=1,\qquad f=2-1.
\]
Their integer polynomials are both $1$. But $\SZ(e)$ asks $1=0$, which is
false, while $\SZ(f)$ asks $2=1$, which is true in this semiring.
The supplied client file constructs the Boolean semiring and checks both
evaluations. This is an explicit negative control on the theorem's scope.

Commutative multiplication also matters. Equality in a commutative
polynomial ring cannot justify $xy=yx$ in a noncommutative matrix semiring.

Conversely, P1 requires no order, characteristic-zero hypothesis, absence of
zero divisors, or multiplicative cancellation. The intended internal-natural
application uses the repository's existing semiring and additive-cancellation
structure. It never identifies a possibly nonstandard model with external
\code{Nat}.

\subsection{Adoption cost and limits}

The supplied core has 72 physical lines, including comments and three axiom
print commands. It is an added interface, whose whole cost belongs in the
accounting. Adopt it when several clients can state their desired balanced
equation and discharge only a polynomial identity, with less dependence on
the recursive split implementation.

Retain the literal expression, polynomial, witness counts, degree
certificates, and all object-model assumptions. The checked core and client
prove the semantic mechanism; they do not certify a full rebuild of the
internal ZFC proof chain. No project-wide line saving is claimed for P1.
